\documentclass[oneside,12pt]{exam}          % ----- exam

%%%%%%%%%%%%%%%%%%%%%%%%%%%%
%%%%% Police et langue %%%%%
%%%%%%%%%%%%%%%%%%%%%%%%%%%%

\usepackage[x11names,table]{xcolor}                           % ----- Couleurs
\usepackage[utf8]{inputenc}                             % ----- UTF-8
\usepackage[french]{babel}                              % ----- French
\usepackage{amsmath,amssymb,mathrsfs}               % ----- Maths
\usepackage{fontawesome5,txfonts}						% ----- Caractères / Symboles
\usepackage{eurosym}                                    % ----- Currencies

\usepackage{helvet}										% ----- Police texte
\renewcommand{\familydefault}{\sfdefault}
\usepackage[T1]{fontenc}

\usepackage{mathpazo}									% ----- Police maths


%%%%%%%%%%%%%%%%
%%%%% TikZ %%%%%
%%%%%%%%%%%%%%%%

\usepackage{pgfplots}                                   % ----- Plots
    \pgfplotsset{compat=1.18}
\usepackage{tikz, tikz-3dplot, tkz-tab , tkz-euclide}                 % ----- TikZ, TikZ 3D, Tableaux

    \usetikzlibrary{math,arrows,arrows.meta}                            % ----- Maths, Vecteurs
    \usetikzlibrary{automata, positioning,patterns.meta}    % ----- Graphs
    \usetikzlibrary{decorations.fractals}                   % ----- Fractals
    \usetikzlibrary{tikzmark}
    \usetikzlibrary{fadings}
    \usetikzlibrary{shapes.geometric}
	\tikzset{>=Stealth}
    
\usepackage{pgf-pie}

    \usetikzlibrary{calc}
    \usetikzlibrary{angles}


\tikzset{arrow style/.style = {\tkzTabDefaultWritingColor,-Stealth,> = \tkzTabDefaultArrowStyle,shorten > = \tkzTabDefaultSep,shorten < = \tkzTabDefaultSep}}		% ----- Modifier flèche tabVar

\tikzset{node style/.style = {inner sep = \tkzTabDefaultSep, outer sep = \tkzTabDefaultSep}}

%%%%%%%%%%%%%%%%%%%%
%%%%% Pratique %%%%%
%%%%%%%%%%%%%%%%%%%%

\usepackage{titlesec}							% ----- Titres et sections
\usepackage{calc}                                       % ----- Calc / Minipage
%\usepackage{fancyhdr}                                   % ----- Header and footer
\usepackage{geometry}                                   % ----- Géométrie de la page
\usepackage{tcolorbox}                                  % ----- Boxes environment
    \tcbuselibrary{skins,listings,listingsutf8}             % ----- All libraries
\usepackage{varwidth}                                   % ----- Adjust box width
\usepackage{fourier-orns}                               % ----- Déco
\usepackage{witharrows}									% ----- Équations avec flèches
\usepackage{array}										% ----- Environnements array / tabular
\usepackage{enumitem}									% ----- Customisation des \enumerate
	\setlist[itemize]{label=$\bullet$}						% ----- Label ● par défaut
\usepackage[c]{esvect}									% ----- Vecteurs
\usepackage{diagbox}									% ----- Découper des cases de \tabular
\usepackage{esvect}
\usepackage{graphicx}									% ----- Insérer des images			% ----- Packages utiles
%%%%%%%%%%%%%%%%%%%%%%%%
%%%%% Mise en page %%%%%
%%%%%%%%%%%%%%%%%%%%%%%%

\usepackage{parskip}									% ----- Parskip
    \setlength{\parskip}{\medskipamount}
    \newcommand{\restoreskip}{\setlength{\parskip}{\medskipamount}}
    
\setlength\parindent{0mm}                               % ----- No indent

\usepackage{setspace}                                   % ----- Interligne
	\setstretch{0.95}                           % Interligne 0.95
\usepackage{makecell}



%%%%%%%%%%%%%%%%%%%%%%%%%%%%%%%
%%%%% Format des sections %%%%%
%%%%%%%%%%%%%%%%%%%%%%%%%%%%%%%

\titlespacing{\section}{0pt}{5mm}{0pt}
\titlespacing{\subsection}{0pt}{2mm}{1mm}


\titleformat{\section}
    {\large\normalfont\bfseries}
    {Exercice \no \thesection\, - \,}
    {0pt}
    {\large\bfseries}

\renewcommand{\thesubsection}{\Alph{subsection}}
\titleformat{\subsection}
    {\normalfont\bfseries}
    {Partie \thesubsection\,: \,}
    {0pt}
    {\bfseries}





		% ----- Préambule feuilles de TD
%%%%%%%%%%%%%%%%%%%%%%%%%%%%%
%%%% Lettres double barre %%%
%%%%%%%%%%%%%%%%%%%%%%%%%%%%%

\newcommand\CC{%
        {\mathbb C}} % ----- Ensemble C

\newcommand\RR{%
        {\mathbb R}} % ----- Ensemble R

\newcommand\QQ{%
        {\mathbb Q}} % ----- Ensemble Q

\newcommand\ZZ{%
        {\mathbb Z}} % ----- Ensemble Z

\newcommand\NN{%
        {\mathbb N}} % ----- Ensemble N

\newcommand\DD{%
        {\mathbb D}} % ----- Ensemble D
        
\newcommand\UU{%
        {\mathbb U}} % ----- Ensemble U

\newcommand\PP{%
        {\mathbb P}} % ----- P probabiliste

\newcommand\EE{%
        {\mathbb E}} % ----- E espérance

\newcommand\VV{%
        {\mathbb V}} % ----- V variance



%%%%%%%%%%%%%%%%%%%%%%%%%%%%%%%%
%%%% Lettres calligraphiques %%%
%%%%%%%%%%%%%%%%%%%%%%%%%%%%%%%%

\newcommand\Ccal{%
        {\mathscr{C}}} % ----- C calligraphique

\newcommand\Dcal{%
        {\mathscr{D}}} % ----- D calligraphique

\newcommand\Pcal{%
        {\mathscr{P}}} % ----- P calligraphique

\newcommand\Rcal{%
        {\mathscr{R}}} % ----- R calligraphique

\newcommand\Fcal{%
        {\mathscr{F}}} % ----- F calligraphique

\newcommand\Scal{%
        {\mathscr{S}}} % ----- S calligraphique

\newcommand\Acal{%
        {\mathscr{A}}} % ----- A calligraphique

\newcommand\Vcal{%
        {\mathscr{V}}} % ----- V calligraphique
        
\newcommand\Tcal{%
        {\mathscr{T}}} % ----- V calligraphique



%%%%%%%%%%%%%%%%%%%%%%%%%%%%%%%%%%%%%
%%% Symboles et/ou calculs usuels %%%
%%%%%%%%%%%%%%%%%%%%%%%%%%%%%%%%%%%%%

%%%%% Analyse %%%%%

\newcommand\limx{%
        {\lim\limits_{x\to +\infty}}\,\,} % ----- Limite x en + infini

\newcommand\limn{%
        {\lim\limits_{n\to +\infty}}\,\,} % ----- Limite n en + infini

\newcommand\limite[1]{%
        {\lim\limits_{\substack{#1}}}\,\,} % ----- Limite X vers Y

\newcommand\nN{%
        {n\in\mathbb N}} % ----- n dans N

\newcommand\xR{%
        {x\in\mathbb R}} % ----- x dans R

\newcommand\zC{%
        {z\in\mathbb C}} % ----- z dans C

\newcommand\somme[2]{%
        {\displaystyle\sum_{#1}^{#2}}}    % ----- Somme

\newcommand\sommek[2]{%
        {\displaystyle\sum_{k=#1}^{#2}}}    % ----- Somme k

\newcommand\inte[2]{%
        {\displaystyle\int_{#1}^{#2}}}   % ----- Intégrale

\newcommand\dd{%
        {\,\,\mathrm{d}}}       % ----- d différentiel

\newcommand\floor[1]{%
        {\lfloor #1 \rfloor}}   % ----- Partie entière inférieure

\newcommand\ceil[1]{%
        {\lceil #1 \rceil}}     % ----- Partie entière supérieure

\newcommand\abs[1]{%
        {\left\lvert #1 \right\rvert}}     % ----- Valeur absolue

\newcommand\pgcd{%
        {\mathrm{pgcd}}}       % ----- pgcd

\newcommand\ppcm{%
        {\mathrm{ppcm}}}       % ----- ppcm



%%%%% Probabilités %%%%%

\newcommand\barre[1]{%
    {\overline{#1}}} % ----- A barre
    
\newcommand\bbigcup{%
    \,\,\bigcup\,\,
    } % ----- Union avec espace 

\newcommand\bbigcap{%
    \,\,\bigcap\,\,
    } % ----- Intersection avec espace


%%%%% Géométrie %%%%%

\newcommand\vcoord[2]{			% ----- Coordonnées d'un vecteur (Nom + coord)
	\renewcommand{\arraystretch}{1.1} \arraycolsep=0.5\arraycolsep
	\vv{#1}\mathbin{=}\begin{pmatrix}#2\end{pmatrix}
} 

\newcommand\Oij{\left(O\,;\,\vv{i}\,;\,\vv{j}\right)} % ----- (O;i;j)

\def\Oijk{\left(O\,;\,\Vec{i}\,;\,\Vec{j}\,;\,\Vec{k}\right)} % ---- (O;i;j;k)

\newcommand\norme[1]{%
        {\left\lVert #1 \right\rVert}}     % ----- Norme

\newcommand\normev[1]{%
        {\left\lVert \vv{#1} \right\rVert}}     % ----- Norme d'un vecteur

\newcommand\scal[2]{%
    {\vv{#1}\cdot\vv{#2}}} % ----- Produit scalaire
    
\newcommand\smeq{{\,=\,}}	% ----- Symbole '=' avec petit espace



%%%%% Algèbre %%%%%

\newcommand\Reel{				% ----- Partie réelle
    {\mathfrak{Re}}} 

\newcommand\Imag{				% ----- Partie imaginaire
    {\mathfrak{Im}}} 

\newcommand\matrice[1]{{ 				% ----- Matrice
    \renewcommand{\arraystretch}{0.7}
    \arraycolsep=0.7\arraycolsep\ensuremath{\begin{pmatrix} #1 \end{pmatrix}}}}

\newcommand\Ouv{(O\,;\,\vv{u}\,;\,\vv{v})} 	% ----- Plan complexe (O;u;v)


%%%%% Flèche SSI longueur adaptative %%%%%

\makeatletter

\newcommand{\ssi}[1]{%
  \@ifnextchar[{\ssi@i{#1}}{\ssi@i{#1}[]}%
}

\def\ssi@i#1[#2]{%
  \@ifnextchar[{\ssi@ii{#1}{#2}}{\ssi@ii{#1}{#2}[]}%
}

\def\ssi@ii#1#2[#3]{%
  \mathrel{%
    \ext@arrow 0099{%
      \Longleftrightarrowfill@
    }%
    {\makebox[#1]{#2}}%
    {\makebox[#1]{#3}}%
  }%
}

\def\Longleftrightarrowfill@{%
  \arrowfill@\Leftarrow\Relbar\Rightarrow
}

\makeatother

                 % ----- Macros mathématiques
%%%%%%%%%%%%%%%%%%%%%%%%%%%%%%%%
%%%%% ----- Packages ----- %%%%%
%%%%%%%%%%%%%%%%%%%%%%%%%%%%%%%%

\usepackage{tcolorbox}                                  % ----- Boxes environment
    \tcbuselibrary{skins,listings,listingsutf8,breakable}             % ----- All libraries
\usepackage[bold=0.7]{xfakebold}


%%%%%%%%%%%%%%%%%%%%%%%%%%%%%%%%%%
%%%%% ----- New colors ----- %%%%%
%%%%%%%%%%%%%%%%%%%%%%%%%%%%%%%%%%

\definecolor{persianblue}{HTML}{221177} % ----- New blues
\definecolor{lightblue}{HTML}{EEFFFF}
\definecolor{softblue}{HTML}{6655BB}

\definecolor{maroon}{HTML}{992222}      % ----- New reds
\definecolor{lightred}{HTML}{FAEEEE}

\definecolor{calgreen}{HTML}{005500}    % ----- New greens
\definecolor{lightgreen}{HTML}{EEFFEE}
\definecolor{softgreen}{HTML}{449944}

\definecolor{brickorange}{HTML}{CC7033} % ----- New oranges
\definecolor{lightorange}{HTML}{FFF2E6}

\definecolor{borderpink}{HTML}{aa00bb} % ----- New pinks
\definecolor{lightpink}{HTML}{ffeeff}

\definecolor{newred}{HTML}{CC0000}     % ----- Colors for geometry
\definecolor{newblue}{HTML}{0000CC}
\definecolor{newgreen}{HTML}{00AA00}
\definecolor{neworange}{HTML}{FFAA00}







%%%%%%%%%%%%%%%%%%%%%%%%%%%%%%%%%%%%%%%%
%%%%% ----- \begin{theoreme} ----- %%%%%
%%%%%%%%%%%%%%%%%%%%%%%%%%%%%%%%%%%%%%%%

\newtcolorbox{theoreme}[2][]{%
    enhanced , breakable , parbox=false ,
    arc=5pt,
    frame style={left color=persianblue, right color=softblue},
    colback=lightblue,
    fonttitle=\bfseries,
    title={\faBook~~ \underline{Théorème} : #2~~\faBook},
    attach boxed title to top left={xshift=5mm , yshift=-2mm},
    boxed title style={
        colback=persianblue,
        colframe=persianblue
    },
    #1
}



%%%%%%%%%%%%%%%%%%%%%%%%%%%%%
%%%%% \begin{propriete} %%%%%
%%%%%%%%%%%%%%%%%%%%%%%%%%%%%

\newtcolorbox{propriete}[2][]{%
    enhanced , breakable , parbox=false ,
    arc=5pt,
    frame style={left color=persianblue, right color=softblue},
    colback=lightblue,
    fonttitle=\bfseries,
    title={\faBook~~ \underline{Propriété} : #2~~\faBook},
    attach boxed title to top left={xshift=5mm,yshift=-2mm},
    boxed title style={
        colback=persianblue,
        colframe=persianblue
    },
    #1
}



%%%%%%%%%%%%%%%%%%%%%%%%%%%%%%%%%%%%%%%%%%
%%%%% ----- \begin{definition} ----- %%%%%
%%%%%%%%%%%%%%%%%%%%%%%%%%%%%%%%%%%%%%%%%%

\newtcolorbox{definition}[2][]{%
    enhanced , breakable , parbox=false ,
    arc=5pt,
    frame style={left color=calgreen, right color=softgreen},
    colback=lightgreen,
    fonttitle=\bfseries,
    title={\faUniversity~~ \underline{Définition} : #2~~\faUniversity},
    attach boxed title to top left={xshift=5mm,yshift=-2mm},
    boxed title style={
        colback=calgreen,
        colframe=calgreen
    },
    #1
}



%%%%%%%%%%%%%%%%%%%%%%%%%%%%%%%%%%%%%%
%%%%% ----- \begin{astuce} ----- %%%%%
%%%%%%%%%%%%%%%%%%%%%%%%%%%%%%%%%%%%%%

\newtcolorbox{astuce}[1][]{%
    enhanced , parbox=false ,
    arc=10pt,
    frame hidden,
    colback=lightred,
    fonttitle=\bfseries,
    borderline={1mm}{0mm}{maroon,dotted},
    title={\faLightbulb~~\underline{Astuce}~~\faLightbulb},
    attach boxed title to top left={xshift=5mm,yshift=-2mm},
    boxed title style={
        colback=maroon,
        colframe=maroon
    },
    #1
}



%%%%%%%%%%%%%%%%%%%%%%%%%%%%%%%%%%%%%%%%
%%%%% ----- \begin{remarque} ----- %%%%%
%%%%%%%%%%%%%%%%%%%%%%%%%%%%%%%%%%%%%%%%

\newtcolorbox{remarque}[1][]{%
    enhanced , parbox=false ,
    arc=10pt,
    frame hidden,
    colback=lightred,
    fonttitle=\bfseries,
    borderline={1mm}{0mm}{maroon,dotted},
    title={\faExclamation~~\underline{Remarque}~~\faExclamation},
    attach boxed title to top left={xshift=5mm,yshift=-2mm},
    boxed title style={
        colback=maroon,
        colframe=maroon
    },
    #1
}



%%%%%%%%%%%%%%%%%%%%%%%%%%%%%%%%%%%%%%%
%%%%% ----- \begin{methode} ----- %%%%%
%%%%%%%%%%%%%%%%%%%%%%%%%%%%%%%%%%%%%%%

\newtcolorbox{methode}[2][]{%
    enhanced , parbox=false ,
    arc=10pt,
    frame hidden,
    colback=lightred,
    fonttitle=\bfseries,
    borderline={1mm}{0mm}{maroon,dotted},
    title={\faCogs~~\underline{Méthode} : #2~~\faCogs},
    attach boxed title to top left={xshift=5mm,yshift=-2mm},
    boxed title style={
        colback=maroon,
        colframe=maroon
    },
    #1
}



%%%%%%%%%%%%%%%%%%%%%%%%%%%%%%%%%%%%%%%
%%%%% ----- \begin{exemple} ----- %%%%%
%%%%%%%%%%%%%%%%%%%%%%%%%%%%%%%%%%%%%%%

\newcommand{\exempledecoration}{%
    \draw[color=brickorange , line width = 2pt , fill=white]
        (frame.south west) ++(0,1pt) -- ++(3cm,0) ++(0,1.5mm) rectangle ++ (3mm,-3mm) ;
}

\newtcolorbox{exemplebox}[1][]{%
	enhanced , breakable , parbox=false ,
    arc=0pt,
    frame hidden,
    borderline west={2pt}{0pt}{brickorange},
    interior style={
        top color=lightorange,
        bottom color=lightorange!15!white
    },
    fonttitle=\bfseries,
    overlay unbroken={\exempledecoration},
    overlay last={\exempledecoration},
    #1
}

\newenvironment{exemple}[1]{%
    \begin{exemplebox}
        \textbf{\color{brickorange}\underline{Exemple} : #1}\\[5pt]
}{%
    \end{exemplebox}
}




%%%%%%%%%%%%%%%%%%%%%%%%%%%%%%%%%%%%%%%%%%%%%
%%%%% ----- \begin{demonstration} ----- %%%%%
%%%%%%%%%%%%%%%%%%%%%%%%%%%%%%%%%%%%%%%%%%%%%

\newcommand{\demodecoration}{%
    \draw[color=borderpink , line width = 2pt , fill=white]
        (frame.south west) ++(0,1pt) -- ++(3cm,0) ++(0,1.5mm) rectangle ++ (3mm,-3mm) ;
}

\newtcolorbox{demobox}[1][]{%
	enhanced , breakable , parbox=false ,
    arc=0pt,
    frame hidden,
    borderline west={2pt}{0pt}{borderpink},
    interior style={
        top color=lightpink,
        bottom color=lightpink!15!white
    },
    fonttitle=\bfseries,
    overlay unbroken={\demodecoration},
    overlay last={\demodecoration},
    #1
}

\newenvironment{demonstration}{%
    \begin{demobox}
        \textbf{\color{borderpink}\underline{Démonstration}}\\[5pt]
}{%
    \end{demobox}
}



%%%%%%%%%%%%%%%%%%%%%%%%%%%%%%%%%
%%%%% ----- \exercice ----- %%%%%
%%%%%%%%%%%%%%%%%%%%%%%%%%%%%%%%%

\newtcolorbox[auto counter , number within=section , number freestyle={\noexpand\arabic{\tcbcounter}} ]{exercicebox}[2][]{% ----- Exercice box
enhanced,
arc=1pt,
colframe=black!80,
colback=black!10,
fonttitle=\bfseries,
hbox,
title={Exercice n° \thetcbcounter \, : #2 },
attach boxed title to top*,
boxed title style = {colback = black!80},
#1}


\newcommand\exercice[2]{% ----- Exercice command
    \begin{exercicebox}[]{#1}
        \begin{varwidth}{\linewidth-30pt} 
         #2
        \end{varwidth}
    \end{exercicebox}
}



%%%%%%%%%%%%%%%%%%%%%%%%%%%%%%%%
%%%%% ----- \calculs ----- %%%%%
%%%%%%%%%%%%%%%%%%%%%%%%%%%%%%%%

\newtcolorbox{calculbox}[1][]{% ----- Calculs box
enhanced,
arc=1pt,
colbacktitle=black!25!white,
coltitle=black,
colback=white,
colframe=black!25!white,
hbox,
title={Aides de calcul},
attach boxed title to top*,
#1}


\newcommand\calculs[1]{% ----- Calculs command
    \begin{calculbox}
        \restoreskip
        #1
    \end{calculbox}
}




%%%%%%%%%%%%%%%%%%%%%%%%%%%%%%%%%%%%%%
%%%%% ----- \begin{rappel} ----- %%%%%
%%%%%%%%%%%%%%%%%%%%%%%%%%%%%%%%%%%%%%

\newtcolorbox{rappel}[2][]{%
    enhanced , breakable , parbox=false ,
    arc=5pt,
    frame hidden,
    colback=lightgreen,
    fonttitle=\bfseries,
    borderline={2pt}{0mm}{calgreen , dashed},
    title={\faUniversity~~ \underline{Rappel} : #2~~\faUniversity},
    attach boxed title to top left={xshift=5mm,yshift=-2mm},
    boxed title style={
        colback=calgreen,
        colframe=calgreen
    },
    #1
}





%%%%%%%%%%%%%%%%%%%%%%%%
%%%%% Mise en page %%%%%
%%%%%%%%%%%%%%%%%%%%%%%%

%%%%% Établissement scolaire %%%%%
\def\etablissement{Lycée Polyvalent Irène et Frédéric Joliot Curie}



%%%%% Adaptation du stretch des environnements 'cases' %%%%%

\makeatletter
\renewcommand*\env@cases[1][1.2]{%
  \let\@ifnextchar\new@ifnextchar
  \left\lbrace
  \def\arraystretch{#1}%
  \array{@{}l@{\quad}l@{}}%
}
\makeatother



%%%%% Ajustement hauteur de \ldots %%%%%
\let\oldldots\ldots
\renewcommand{\ldots}{%
	\parbox[b][15pt]{15pt}{\oldldots}
}


%%%%% Colonne centrée de taille ajustable %%%%%
\newcolumntype{P}[1]{>{\centering\arraybackslash}m{#1}<{\vspace{0pt}}}



%%%%%%%%%%%%%%%%%%%%%%%%%%%%%%%%%%%%%%
%%% Boîte "codebox" pour les algos %%%
%%%%%%%%%%%%%%%%%%%%%%%%%%%%%%%%%%%%%%

\newcommand{\CodeNumber}[1]{%
    \makebox[0pt][c]{#1\hspace{4mm}}%
}

%%%%% Style listings %%%%%
\lstdefinestyle{moncode}{
    basicstyle = \ttfamily,
    numbers = left,
    numberstyle = \footnotesize\CodeNumber,
    numbersep = 4mm,
    numberblanklines = true,
    breaklines = true,
    columns = fullflexible,
    keepspaces = true,
    showstringspaces = false,
    tabsize = 4,
    frame = none
}

%%%%% Environnement codebox %%%%%
\newtcblisting{codebox}[1][]{
	hbox , 
    listing only , 
	listing engine = listings ,
    enhanced ,
	overlay = {
	    \draw[newred , line width = 1pt]
		([xshift = 7mm,yshift=-2mm]frame.north west)
		--
		([xshift = 7mm,yshift=3mm]frame.south west);
		},
    listing options = {
    	style = moncode ,
    	xleftmargin = 0mm
    	},
    colframe=newred , colback=red!15 ,
    arc = 2pt , boxrule = 1.5pt,
	left = 8mm , right = 3mm , top = 1mm , bottom = 1mm ,
    #1
}



%%%%%%%%%%%%%%%%%%%%%%%%%%%%%%%%%%%%%%%%%%%%%%%%%%%%%%%%%
%%%%% Commande \lignes{} pour remplir avec \dotfill %%%%%
%%%%%%%%%%%%%%%%%%%%%%%%%%%%%%%%%%%%%%%%%%%%%%%%%%%%%%%%%

\makeatletter
\newcommand{\lignes}[1]{%
  \begingroup
  \fontdimen3\font=8pt
  \dotfill
  \count@=1
  \loop
    \ifnum\count@<#1
      \par
      \vspace{8pt}
      \noindent\dotfill
      \advance\count@ by 1
  \repeat
  \endgroup
}
\makeatother 		   	% ----- Macros de mise en page

\geometry{left=30pt, right=30pt, top=75pt, bottom=55pt} % ----- Marges

\begin{document}


\firstpageheader{Prénom : ........................................\\[3mm] Nom : ............................................. }{}{Classe : ............... \\[3mm] Date : ...............}
\footer{}{\footnotesize \etablissement}{}

\hrulefill\raisebox{-1.9pt}{\quad\decofourleft~\decosix~\decofourright \quad}\hrulefill

\begin{center}
	\textbf{\large Devoir surveillé \no 2 (55 min)}
\end{center}
\vspace{-5pt}
Vous êtes invité\textperiodcentered e\textperiodcentered s à faire figurer sur votre copie toute trace de recherche, même incomplète ou non fructueuse, que vous aurez développée. La qualité de votre rédaction, la clarté et la précision de vos raisonnements seront prises en compte dans l’appréciation de votre copie.

Les élèves disposant d'un P.A.P. n'ont pas à traiter les questions marquées du symbole $\bigstar$

\vspace{-8pt}\hrulefill\raisebox{-1.9pt}{\quad\decofourleft~\decosix~\decofourright \quad}\hrulefill
\vspace{5pt}


%%%%%%%%%%%%%%%%%%%%%%%%%%%%%%%%%%%%%%%%%%%%%%%%%%%%%%%%%%%%%%%%%%%%%%%%%%%%%%%%%%%%%%%%%%%%%%%%%%%%%%%%%%%%%
%%%%%%%%%%%%%%%%%%%%%%%%%%%%%%%%%%%%%%%%%%%%%%%%%%%%%%%%%%%%%%%%%%%%%%%%%%%%%%%%%%%%%%%%%%%%%%%%%%%%%%%%%%%%%
%%%%%%%%%%%%%%%%%%%%%%%%%%%%%%%%%%%%%%%%%%%%%%%%%%%%%%%%%%%%%%%%%%%%%%%%%%%%%%%%%%%%%%%%%%%%%%%%%%%%%%%%%%%%%

\section{Questions de cours}
\begin{enumerate}[itemsep = 0pt]
	\item Qu'est-ce qu'une suite arithmétique ?
	\item Donner la formule permettant d'exprimer $u_n$ en fonction de $n$ si son premier terme est $u_1$
	\item $\bigstar$ À quoi correspond le symbole $\Sigma$ ? Donner son nom et sa fonction.
\end{enumerate}



\section{Généralités sur les suites}
\begin{enumerate}[itemsep = 0pt]
	\item On considère la suite $(u_n)$ définie par $u_n = n^2 - n + 1$ 
	
	Donner les $5$ premiers termes de la suite $(u_n)$
	
	\item $\bigstar$ On considère la suite $(v_n)$ définie par $\begin{cases} v_0 = 1 \\ v_{n+1} = 2v_n +1 \end{cases}$ 
	
	Donner les $5$ premiers termes de la suite $(v_n)$
	
	\item Un nouvel opérateur téléphonique ouvre et réussit à démarcher 15\,000 abonnés le premier mois. Chaque mois suivant, cet opérateur perd $10\%$ de ses abonnés mais réussit à en démarcher 5\,000 nouveaux. \\
	Donner une définition de $(w_n)$ la suite donnant le nombre d'abonnés de l'opérateur téléphonique après $n$ mois.
\end{enumerate}



\section{À petit pas}
On s'intéresse à une échelle dont le 8e barreau se situe à 3,45 m de hauteur, et dont le 12e barreau se situe à 5,25 m.

On note $u_n$ la hauteur du $n$-ième barreau de l'échelle exprimée en mètres. On admettra que $(u_n)$ est une suite arithmétique, autrement dit l'écart entre les barreaux est constant.
\begin{enumerate}[itemsep = 0pt]
	\item Traduire les valeurs de l'énoncé en donnant la valeur de deux termes de la suite $(u_n)$
	\item Montrer que la raison de la suite $(u_n)$ est $0,45$ en justifiant par une formule du cours.
	\item Retrouver la hauteur du premier barreau de l'échelle.
	\item On admettra que l'expression de la suite $(u_n)$ est $u_n = 0,45n - 0,15$
	\begin{enumerate}
		\item Donner la hauteur du 5e barreau de l'échelle.
		\item $\bigstar$ Sachant que le dernier barreau de l'échelle se trouve à une hauteur de 8,85m, combien de barreau l'échelle comporte-t-elle ?
	\end{enumerate}
\end{enumerate}



%%%%%%%%%%%%%%%%%%%%%%%%%%%%%%%%%%%%%%%%%%%%%%%%%
\newpage
\newgeometry{left=30pt,right=30pt, top=40pt, bottom=55pt}
%%%%%%%%%%%%%%%%%%%%%%%%%%%%%%%%%%%%%%%%%%%%%%%%%



\section{Autour de $\Sigma$}
\begin{enumerate}[itemsep = 0pt]
	\item Écrire chacune des sommes suivantes en utilisant le symbole $\Sigma$ sans essayer de les calculer :

	\begin{tabular}{p{0.3\textwidth} p{0.3\textwidth} p{0.3\textwidth}}
		$A = 4^3 + 5^3 + 6^3 + \ldots + 21^3$ &
		$\bigstar$ $B = 5^5 + 6^6 + 7^7 + \ldots + 18^{18}$ &
		$C = \dfrac{7}{8} + \dfrac{7}{9} + \dfrac{7}{10} + \ldots + \dfrac{7}{77}$
	\end{tabular}
	
	\item On considère $(u_n)$ une suite arithmétique de premier terme $u_0 = 1$ et de raison $r = 3$
	\begin{enumerate}
		\item Exprimer $(u_n)$ en fonction de $n$
		\item Calculer la somme $S = \sommek{7}{15} u_k$
	\end{enumerate}
\end{enumerate}



\section{Piqure de rappel}
Dans une urne on place 40 boules numérotées de $1$ à $40$. On considère les évènements suivants :

$\bullet$ A : \og La boule tirée porte un numéro pair \fg \\
$\bullet$ B : \og La boule tirée porte un numéro qui est un multiple de 3\fg \\
$\bullet$ C : \og La boule tirée porte un numéro supérieur ou égal à $25$ \fg
\begin{enumerate}[itemsep = 0pt]
	\item Déterminer $\PP(A)$ , $\PP(B)$ et $\PP(C)$
	\item Décrire explicitement l'évènement $A \cap B$
	\item En justifiant, déterminer la valeur de $\PP(B \cap C)$
\end{enumerate}



%%%%%%%%%%%%%%%%%%%%%%%%%%%%%%%%%%%%%%%%%%%%%%%%%
%%%%%%%%%%%%%%%%%% CORRECTION %%%%%%%%%%%%%%%%%%%
%%%%%%%%%%%%%%%%%%%%%%%%%%%%%%%%%%%%%%%%%%%%%%%%%
\newpage
\setcounter{section}{0}
\color{newblue}
%%%%%%%%%%%%%%%%%%%%%%%%%%%%%%%%%%%%%%%%%%%%%%%%%



\section{Questions de cours}
\begin{enumerate}[itemsep = 0pt]
	\item Une suite est arithmétique si la différence entre deux de ses termes consécutifs est constante.
	\item $u_n = u_1 + (n-1) \times r$
	\item $\bigstar$ Le symbole $\Sigma$ se nomme "sigma", il permet de décrire une somme.
\end{enumerate}



\section{Généralités sur les suites}
\begin{enumerate}[itemsep = 0pt]
	\item On considère la suite $(u_n)$ définie par $u_n = n^2 - n + 1$ 
	
	$u_0 = 0^2 - 0 + 1 = 1$ \qquad $u_1 = 1^2 - 1 + 1 = 1$ \qquad $u_2 = 2^2 - 2 + 1 = 3$ 
	
	$u_3 = 3^2 - 3 + 1 = 7$ \qquad $u_4 = 4^2 - 4 + 1 = 13$
	
	\item $\bigstar$ On considère la suite $(v_n)$ définie par $\begin{cases} v_0 = 1 \\ v_{n+1} = 2v_n +1 \end{cases}$ 
	
	$v_0 = 1$ \qquad $v_1 = 2v_0 + 1 = 3$ \qquad $v_2 = 2v_1 + 1 = 7$ \qquad
	$v_3 = 2v_2 + 1 = 15$ \qquad $v_4 = 2v_3 + 1 = 31$
	
	\item $\begin{cases}w_0 = 15\,000 \\ w_{n+1} = 0,9w_n + 5000 \end{cases}$
\end{enumerate}



\section{À petit pas}
On s'intéresse à une échelle dont le $8$-ième barreau se situe à 3,45m de hauteur, et dont le $12$-ième barreau se situe à 5,25m.

On note $u_n$ la hauteur du $n$-ième barreau de l'échelle exprimée en mètres. On admettra que $(u_n)$ est une suite arithmétique, autrement dit l'écart entre les barreaux est constant.
\begin{enumerate}[itemsep = 0pt]
	\item D'après l'énoncé : $u_8 = 3,45$ et $u_{12} = 5,25$
	\item $r = \dfrac{u_{12}-u_8}{12-8} = \dfrac{5,25 - 3,45}{4} = \dfrac{1,8}{4} = 0,45$
	\item Pour retrouver le premier terme de la suite il faut retirer 7 fois la raison à $u_8$
	
	On a donc $u_1 = 3,45 - 7 \times 0,45 = 0,3$ , le premier barreau de l'échelle se situe à 30cm de hauteur
	\item On admettra que l'expression de la suite $(u_n)$ est $u_n = 0,45n - 0,15$
	\begin{enumerate}
		\item $u_5 = 0,4 \times 5 - 0,15 = 2,1$ , le $5$-ième barreau de l'échelle se situe à 2,10m de hauteur.
		\item $\bigstar$ On sait que $u_n = 0,45n - 0,15$ , et on cherche la valeur de $n$ pour laquelle $u_n = 8,85$
		
		Or $0,45n - 0,15 = 8,85 \iff 0,45n = 9 \iff n = \dfrac{9}{0,45} \iff n = 20$
		
		On peut conclure que cette échelle comporte 20 barreaux.
	\end{enumerate}
\end{enumerate}



\section{Autour de $\Sigma$}
\begin{enumerate}[itemsep = 0pt]
	\item Écrire chacune des sommes suivantes en utilisant le symbole $\Sigma$ sans essayer de les calculer :

		$A = 4^3 + 5^3 + 6^3 + \ldots + 21^3 = \sommek{4}{21} k^3$ 
		
		$\bigstar$ $B = 5^5 + 6^6 + 7^7 + \ldots + 18^{18} = \sommek{5}{18} k^k$ 
		
		$C = \dfrac{7}{8} + \dfrac{7}{9} + \dfrac{7}{10} + \ldots + \dfrac{7}{77} = \sommek{8}{77} \dfrac{7}{k}$
	
	\item On considère $(u_n)$ une suite arithmétique de premier terme $u_0 = 1$ et de raison $r = 3$
	\begin{enumerate}
		\item $u_n = u_0 + n\times r = 1 + 3n$
		\item On applique la formule $\sommek{p}{m} u_k = (m-p+1) \times \dfrac{u_m + u_p}{2}$
		
		On commencer par calculer $u_7 = 1 + 7\times 3 = 22$ et $u_15 = 1+ 15 \times 3 = 46$
		
		Finalement : $S = \sommek{7}{15} u_k = (15-7+1) \times \dfrac{u_{15} + u_7}{2} = 9 \times \dfrac{22 + 46}{2} = 9 \times 34 = 306$
	\end{enumerate}
\end{enumerate}



\section{Piqure de rappel}
Dans une urne on place 40 boules numérotées de $1$ à $40$. On considère les évènements suivants :

$\bullet$ A : \og La boule tirée porte un numéro pair \fg \\
$\bullet$ B : \og La boule tirée porte un numéro qui est un multiple de 3\fg \\
$\bullet$ C : \og La boule tirée porte un numéro supérieur ou égal à $25$ \fg
\begin{enumerate}[itemsep = 0pt]
	\item $\PP(A) = \dfrac{1}{2}$ \quad $\PP(B) = \dfrac{13}{40}$ \quad $\PP(C) = \dfrac{16}{40} = \dfrac{2}{5}$
	\item $A \cap B = \{6 ~;~ 12 ~;~ 18 ~;~ 24 ~;~ 30 ~;~ 36\}$
	\item $A \cap C = \{26 ~;~ 28 ~;~ 30 ~;~ 32 ~;~ 34 ~;~ 36 ~;~ 38 ~;~ 40\}$
	
	Donc $\PP(A \cap B) = \dfrac{8}{40} = \dfrac{1}{5}$
\end{enumerate}


\end{document}